\section{Hypotheses for the global connection problem}
\label{sec:supp-connection-details}

This section gives a precise functional-analytic form of the additional
hypotheses used in
\cref{sec:heteroclinic-connection,prop:connection-comparison}.  We work
with the globally modified equation defined by
\cref{eq:rw-voltage-rfde,eq:rw-modified-recovery}; the function \(g\) is
fixed throughout.  Uniformity in \(N\) and in the parameter box from
\cref{cor:rw-conditional-connection} is part of each stated hypothesis.
Let \(\mathcal F_p^g\) denote this equation's fold-time RFDE vector field,
where \(p=(\nu,\eta)\), and put
\[
 \sigma(r)=(r-\delta^2\nu)g(r),
 \qquad \sigma_\pm'=\sigma'(\pm r_A).
\]

Except in the characteristic-root calculation at the equilibria, we use
fold time \(s=\delta t\), so the delay translations are
\(s\mapsto s-\theta_k\).  The spectral parameter in that calculation is
explicitly denoted \(\lambda_t\) and belongs to the original time \(t\).
If \(\Sigma\) is a regular transverse section, then
\(\mathcal M_\Sigma^2\) denotes the solution-manifold section from
\cref{eq:C2-solution-manifold-section}, and
\(\mathcal P_{\Sigma_1\Sigma_2,p}\) denotes the forward transition map
between two such sections.  The transition map is used only where the
crossing speed has
a common nonzero lower bound.

Fix \(0<r_L<r_{\rm loc}\) and write
\begin{equation}
 \Sigma_L=\{\phi:r(\phi)=-r_L\}.
 \label{eq:supp-negative-regular-section}
\end{equation}
A \emph{negative reference branch} below is a \(C^2\) orbit segment
\(Z_p^-(s)=(v_p^-(s),w_p^-(s))\) of the fixed modified RFDE, parameterized in
fold time, for which
\[
 r_p(s):=\pi_N^T\{v_p^-(s)-\mathbf1\}<0,
 \qquad r_p'(s)<0,
\]
and whose forward continuation converges to \(Z_N^-\).  The branch used in
the proofs is the one obtained from the stable manifold in
\cref{hyp:supp-outer-manifolds}(iii).  On an attracting positive reference
segment we use the same notation \(r_p(s)\), now with \(r_p(s)>0\) and
\(r_p'(s)<0\).  Let
\begin{equation}
 \mathcal J_N^{\rm fv}(r)
 =D(P_N-\Id)-\diag\!\left(
 2c_N\circ\{V_{0,N}(r)-\mathbf1\}
 +\beta\{V_{0,N}(r)-\mathbf1\}^{\circ2}\right)
 \label{eq:supp-frozen-voltage-jacobian}
\end{equation}
and denote by \(\lambda_{c,N}(r)\) its simple real eigenvalue near zero; it
is positive for \(-r_A<r<0\) and negative for \(0<r<r_A\).  On either
regular branch set
\begin{equation}
 a_p(s)=\left\{\abs{\lambda_{c,N}(r_p(s))}^{-1}
                    +2(D\gamma)^{-1}\right\}^{-1}.
 \label{eq:supp-capped-voltage-rate}
\end{equation}
Thus \(a_p(s)\) is comparable to
\(\min\{\abs{\lambda_{c,N}(r_p(s))},D\gamma\}\), and hence to
\(\abs{r_p(s)}\) on either regular branch.  Choose a fixed
\(c_\sharp\in(0,1/4)\) and
set
\begin{equation}
 \kappa_u(s)=c_\sharp
 \left\{\frac{\delta}{a_p(s)}
       +\frac{\theta_m}{\log(1/\delta)}\right\}^{-1},
 \qquad
 \mathcal A_u(s,t)=\int_t^s\kappa_u(\zeta)\,\dd\zeta,
 \label{eq:supp-dichotomy-action}
\end{equation}
For a voltage history \(\phi\), and for a full history
\((\phi,\omega)\) in which \(\omega\) is the present recovery component,
define
\begin{align}
 \norm{\phi}_{p,s,\sharp}^{v}
 &=\sup_{-\theta_m\le\theta\le0}
   e^{-2\mathcal A_u(s,s+\theta)}\norm{\phi(\theta)}_N,\notag\\
 \norm{(\phi,\omega)}_{p,s,\sharp}
 &=\max\left\{\norm{\phi}_{p,s,\sharp}^{v},
       \frac{\abs\omega}{\abs{w_{0,N}'(r_p(s))}}\right\}.
 \label{eq:supp-fading-history-norm}
\end{align}
This scaling separates the slow tangential direction from the voltage
dichotomy.

For a map depending on \(p=(\nu,\eta)\), write
\[
 \norm{\mathbf D_p^1F}
 =\max\{\norm F,\norm{\partial_\nu F},\norm{D_\eta F}\}.
\]
The last norm is the full operator norm on
\(\mathcal T_N\).
Every statement below that uses a global orbit segment or an invariant
manifold associated with \(Z_N^\pm\) is conditional on
\(\mathrm{(H_{\rm glob})}\) in
\cref{hyp:supp-outer-manifolds}.  The local spectral calculation verifies
only item~\textup{(i)} of that hypothesis.
Throughout the estimates, \(M\ge0\) denotes a fixed exponent larger than
the finitely many losses caused by history traces, transition maps, and one
parameter derivative.  For each prescribed algebraic order, \(p_0\) in
\cref{eq:supp-central-cuts} is then chosen to satisfy
\cref{eq:supp-central-action-choice}.

\subsection{Invariant manifolds at the two equilibria}

\begin{assumption}[Global invariant-manifold and comparison hypotheses]
\label[assumption]{hyp:supp-outer-manifolds}
Fix \(r_A\), \(\delta_0\), and the parameter box in
\cref{cor:rw-conditional-connection}.  We say that
\(\mathrm{(H_{\rm glob})}\) holds when the following objects and estimates are
available uniformly in \(N\) and in that parameter box.

\begin{enumerate}[label=\textup{(\roman*)}]
\item The characteristic roots at \(Z_N^\pm\) are separated from the
imaginary axis.  At \(Z_N^+\) the only unstable root is the simple slow
root
\begin{equation}
 \lambda_{s,\pm}^{t}
 =\delta^2\frac{\sigma_\pm'}{w_{0,N}'(\pm r_A)}+O(\delta^4),
 \label{eq:supp-outer-slow-roots}
\end{equation}
where the superscript indicates the original-time spectral parameter.  At
\(Z_N^-\) this root is stable, and the only unstable root is the simple
voltage root.

\item The local unstable manifold of \(Z_N^+\) is \(C^2\) in the state and
\(C^1\) in \(p\), on a ball of radius \(c_A\delta^4\) in the scaled history
chart.  Its inward component contains a unique orbit defined
for all negative time and converging to \(Z_N^+\) as \(s\to-\infty\).
Forward continuation gives a unique transverse intersection point
\(A_p^0\in\mathcal M_{\Sigma_0}^2\).  On each fixed positive regular
subinterval, the orbit and its first parameter derivatives are bounded in
the norm \cref{eq:supp-fading-history-norm}.

\item The local stable manifold of \(Z_N^-\) is \(C^2\) in the state and
\(C^1\) in \(p\), on a ball of radius \(c_A\delta^4\) in the scaled history
chart.  On a fixed negative regular section \(\Sigma_L\), the histories whose
forward orbits converge to \(Z_N^-\) form a codimension-one \(C^2\)
stable manifold.  There is a \(C^1\) defining function
\(\mathfrak m_{L,p}\) and a normalized vector in the unstable tangent space
\(\jmath_{L,p}^{u,\infty}\) such that, at a reference point
\(\bar Z_{-,L}\) of this manifold,
\begin{equation}
 \Lambda_{L,p}^{\infty}
 :=D\mathfrak m_{L,p}(\bar Z_{-,L}),
 \qquad
 \Lambda_{L,p}^{\infty}\jmath_{L,p}^{u,\infty}=1.
 \label{eq:supp-terminal-normalization}
\end{equation}
The orbit through \(\bar Z_{-,L}\) is defined on a negative half-line and
converges to \(Z_N^-\).

\item Replace the asymptotic condition at \(Z_N^-\) by a boundary condition
on a truncated interval whose endpoint lies a distance \(L_\delta\) beyond
\(\Sigma_L\) along the negative branch.  The integrated unstable
rate
satisfies
\begin{equation}
 \mathcal A_u(L_\delta,0)
 \ge c\frac{\log(1/\delta)}{\delta}.
 \label{eq:supp-integrated-unstable-rate}
\end{equation}
Consequently, replacing the stable-manifold condition at infinity by a
boundary functional at the truncated endpoint, with polynomially bounded
first derivatives, changes the
pulled-back defining covector, after normalization on the same transported
unstable tangent vector, by at most
\begin{equation}
 C\delta^{-M}
 \exp\left\{-c\frac{\log(1/\delta)}{\delta}\right\}.
 \label{eq:supp-terminal-decay}
\end{equation}
The same estimate holds for the defining-function value and after one
state or parameter derivative.

\item Let \(A_p^{\rm fin}\) be the attracting trace used to define
\(D_N^{\rm fin}\) in \cref{prop:finite-comparison}, intersected with
\(\Sigma_0\) after both forward orbits have entered the region where the
two equations agree.  For every \(q>0\), the endpoint in the
finite-interval construction may be placed at a sufficiently large value of
\(S_\delta\) that the Gaussian tail estimate gives
\begin{equation}
 \norm{\mathbf D_p^1(A_p^0-A_p^{\rm fin})}
 \le C_q\delta^q.
 \label{eq:supp-positive-history-comparison}
\end{equation}
\end{enumerate}
\end{assumption}

Only the spectral condition in item~\textup{(i)} is verified below from the
displayed network assumptions.  Items~\textup{(ii)}--\textup{(v)} are
additional hypotheses; local hyperbolicity alone does not supply the global
orbit segments, their intersection, or the comparison estimates.

\begin{lemma}[Spectral condition at the two equilibria]
\label[lemma]{lem:supp-outer-spectral}
Under the assumptions of \cref{cor:rw-conditional-connection}, after
\(r_A\) and \(\delta_0\) are decreased, item~\textup{(i)} of
\(\mathrm{(H_{\rm glob})}\) holds.
\end{lemma}

\begin{proof}
For this spectral calculation only, derivatives and characteristic roots are
taken in the original time \(t\).  Put
\begin{equation}
 J_\pm=D(P_N-\Id)-\diag\!\left(
 2c_N\circ\{V_{0,N}(\pm r_A)-\mathbf1\}
 +\beta\{V_{0,N}(\pm r_A)-\mathbf1\}^{\circ2}\right)
 =\mathcal J_N^{\rm fv}(\pm r_A)
 \label{eq:supp-outer-voltage-jacobian}
\end{equation}
and
\begin{equation}
 A_\pm(\lambda_t)
 =J_\pm+\delta^2K\sum_{k=0}^m(B_{k,N}+\Delta B_{k,N})
       \{1-e^{-\lambda_t\theta_k/\delta}\}.
 \label{eq:supp-original-time-characteristic-block}
\end{equation}
The Dobrushin inverse on \(E_N\), followed by the two invariant-graph
equations for the splitting
\(\R^N=\operatorname{span}\{\mathbf1\}\oplus E_N\), gives one simple real
eigenvalue
\[
 \lambda_{c,N}(r)=-2\alpha r+O(r^2)
\]
of \(\mathcal J_N^{\rm fv}(r)\); the complementary semigroup decays at rate
at least \(D\gamma/2\).  Thus, after \(r_A\) is fixed sufficiently small,
\(J_+\) has no unstable eigenvalue, \(J_-\) has exactly the simple unstable
eigenvalue \(\lambda_{c,N}(-r_A)\), and the resolvents on the imaginary axis
are bounded uniformly in \(N\).  On \(\Re\lambda_t\ge0\), the delay term in
\cref{eq:supp-original-time-characteristic-block} is an
\(O(\delta^2)\) perturbation in the layer norm.  A Neumann-series argument
therefore gives the same imaginary-axis resolvent bound for
\(\lambda_t\Id-A_\pm(\lambda_t)\) when \(\delta\) is small; explicitly,
\begin{equation}
 \sup_{\omega\in\R}
 \norm{(i\omega\Id-A_\pm(i\omega))^{-1}}
 \le C_A,
 \label{eq:supp-outer-imaginary-resolvent}
\end{equation}
where \(C_A\) is independent of \(N\), \(\delta\), and the admissible
constrained perturbation (with the fixed choice of \(r_A\)).

The modal equations are
\[
 (\lambda_t\Id-A_\pm(\lambda_t))u=-\omega\mathbf1,
 \qquad
 \lambda_t\omega=\delta^2\sigma_\pm'\pi_N^Tu.
\]
Eliminating \(u\) near zero gives
\[
 \Delta_\pm(\lambda_t)
 =\lambda_t+\delta^2\sigma_\pm'\pi_N^T
   (\lambda_t\Id-A_\pm(\lambda_t))^{-1}\mathbf1.
\]
Differentiating \cref{eq:rw-critical-curve-equation} yields
\[
 \pi_N^T(-J_\pm)^{-1}\mathbf1
 =-\frac1{w_{0,N}'(\pm r_A)}.
\]
It follows that
\[
 \Delta_\pm(\lambda_t)
 =\lambda_t-\delta^2
   \frac{\sigma_\pm'}{w_{0,N}'(\pm r_A)}+O(\delta^4),
 \qquad
 \Delta_\pm'(\lambda_t)=1+O(\delta^2).
\]
This proves the simple-root expansion in
\cref{eq:supp-outer-slow-roots}.  To count the other roots, replace
\(\delta^2\) in the second modal equation by
\(\varepsilon\in[0,\delta^2]\).  The same imaginary-axis resolvent
excludes a crossing outside a small neighborhood of zero.  Inside that
neighborhood the real part of
\(\varepsilon\sigma_\pm'/w_{0,N}'(\pm r_A)\) excludes a crossing away
from the simple root.  The argument principle therefore preserves the
remaining right-half-plane count.  The contour is taken in the original-time
\(\lambda_t\)-plane; the uniform voltage resolvent excludes imaginary-axis
zeros along the homotopy, and standard retarded-equation root bounds close
the contour in a common right half-plane.  The signs of \(w_{0,N}'\),
\(g'\), and \(\sigma_\pm'\) give the indices in (i).  All subsequent orbit
and dichotomy estimates return to fold time \(s=\delta t\).
\end{proof}

\begin{remark}[Sufficient quantitative data for the remaining hypotheses]
The following conditions are sufficient for the later comparison argument;
they are not derived here from the local spectral calculation.  For
item~\textup{(ii)}, it is sufficient to verify that the local
Lyapunov--Perron map, in \cref{eq:supp-fading-history-norm} and with the slow
coordinate divided by \(\delta^4\), is a uniform contraction on the stated
ball; that its inward one-dimensional branch exists for all negative time;
and that its forward continuation reaches \(\Sigma_0\) with a uniformly
nonzero crossing speed.  The RFDE invariant-manifold theorem and the
implicit-function theorem for transverse intersections then give the regularity asserted in
item~\textup{(ii)}.  This verification uses only forward continuation from
the local unstable manifold.

For item~\textup{(iii)}, one checkable sufficient condition on a fixed
negative regular interval and its delay enlargement is as follows.  Require
an approximate orbit segment of the form
\begin{equation}
 \Gamma_{-,p}^{[10]}(s)=\sum_{j=0}^{9}\delta^j\Gamma_{j,p}(s)
 \label{eq:supp-negative-slow-expansion}
\end{equation}
with \(\Gamma_{0,p}\) on the critical curve
\cref{eq:rw-critical-curve}.  The coefficient equations obtained by
substitution into the fold-time RFDE must be solvable successively.  At each
order, it is enough that the voltage coefficient and scalar speed correction
solve the bordered equation associated with
\begin{equation}
 \mathscr B_N(r):E_N\times\R\longrightarrow\R^N,
 \qquad
 \mathscr B_N(r)(Z,Q)=\mathcal J_N^{\rm fv}(r)Z-Q\mathbf1.
 \label{eq:supp-negative-bordered-map}
\end{equation}
At \(r=0\), projection by \(\pi_N^T\) determines \(Q\), while projection by
\(P_{\perp,N}\) determines \(Z\) through \(A_N^{-1}\); the Dobrushin bound
therefore gives a uniform inverse for \(\mathscr B_N(0)\).  A sufficient
regular-interval condition is
\[
 \sup_{N,p,\,r\in[-r_A,-r_L]}
 \norm{\mathscr B_N(r)^{-1}}<\infty,
 \qquad
 \inf_{N,p,\,r\in[-r_A,-r_L]}\abs{w_{0,N}'(r)}>0.
\]
Together with the fixed delay translations and the uniform layer bounds,
these inequalities make the coefficient recursion and its first
\(p\)-derivative uniform.  Write \(\Gamma_-=\Gamma_{-,p}^{[10]}\) and define
\begin{align}
 R_{\Gamma_-}(s)
 &=\Gamma_-'(s)-\mathcal F_p^g((\Gamma_-)_s),\notag\\
 \mathcal N_{\Gamma_-}(z)
 &=\mathcal F_p^g((\Gamma_-+z)_s)
   -\mathcal F_p^g((\Gamma_-)_s)
   -D\mathcal F_p^g((\Gamma_-)_s)z_s.
 \label{eq:supp-negative-residual-definition}
\end{align}
The required residual estimate is
\begin{equation}
 \norm{R_{\Gamma_-}}+\norm{D_pR_{\Gamma_-}}\le C\delta^9.
 \label{eq:supp-negative-residual}
\end{equation}
In addition, require an exponential dichotomy on the negative half-line with
a Green operator that propagates the stable complement forward and evaluates
the unstable coefficient by
\begin{equation}
 (\mathbb G_-f)_u(s)
 =-\int_s^\infty U_u(s,\zeta)P_u(\zeta)f(\zeta)\,\dd\zeta.
 \label{eq:supp-half-line-green}
\end{equation}
and satisfies, in the scaled range and history norms,
\[
 \norm{\mathbb G_-}\le C\delta^{-2}.
\]
Finally require that, on a ball of radius \(c\delta^6\), the map
\[
 z\longmapsto
 \mathbb G_-\{-R_{\Gamma_-}+\mathcal N_{\Gamma_-}(z)\}
\]
maps the ball into itself and has Lipschitz constant at most \(1/3\), with the
same bounds after one parameter derivative.  These conditions imply the
stable manifold and its \(C^1\) parameter dependence through the standard
Lyapunov--Perron theorem; varying the stable entry coordinate supplies the
local graph, and \cref{eq:supp-terminal-normalization} fixes its scalar
normalization.  They are a verification route for item~\textup{(iii)}, not a
consequence of item~\textup{(i)} alone.

For item~\textup{(iv)}, it is sufficient to verify the action lower bound
\cref{eq:supp-integrated-unstable-rate}, a polynomial bound for the endpoint
functional and its first derivatives, and the corresponding roughness bound
for the dichotomy.  The difference between the infinite unstable integral in
\cref{eq:supp-half-line-green} and its truncated counterpart is then confined
to the omitted tail.  Applying the dichotomy estimate to this tail, followed
by the mean-value identity for the two Lyapunov--Perron fixed points, yields
\cref{eq:supp-terminal-decay}; differentiation yields the stated derivative
estimate.

For item~\textup{(v)}, it is sufficient that the two attracting histories,
once they satisfy the same RFDE and use the same time parametrization, obey a forward
variation-of-constants estimate with polynomially bounded initial
difference, attracting factor
\(e^{-c\log(1/\delta)/\delta}\), and fold-passage tail
\[
 C\delta^{-M}e^{-S_\delta^2/2+cS_\delta}
       \langle S_\delta\rangle^M.
\]
For each fixed \(q>0\), choosing the exponent in \(S_\delta\) sufficiently
large makes both terms \(O(\delta^q)\).  The differentiated equations have
only a finite
additional power loss, which is absorbed by the same choice.
\end{remark}

\subsection{A bounded solution operator for the normal variational equation}

\begin{lemma}[Normal variational equation in the fold chart]
\label[lemma]{lem:supp-central-boundary-inverse}
Let \(\mathcal I_p\) be the phase-space embedding associated with the invariant graph in
\cref{prop:fold-invariant-histories}, and let \(e_{\rm pl}\) denote evaluation
at \(\theta=0\) of the two fold variables.  Let
\(I=[s_-,s_+]\subset[-S_\delta,S_\delta]\) be a closed fold-chart interval such
that its delay enlargement
\(I^-=[s_--\theta_m,s_+]\) remains in the region where the original and
modified equations agree.  Along an orbit \(u_p:I\to\R^2\) on the invariant
graph, the map
\begin{equation}
 \phi\longmapsto
 \bigl(e_{\rm pl}\phi,\,
       \phi-\mathcal I_p(e_{\rm pl}\phi)\bigr)=:(u,\omega)
 \label{eq:supp-central-shear}
\end{equation}
is a local \(C^2\)-in-state and \(C^1\)-in-\(p\) diffeomorphism.  Put
\[
 \mathscr H_{\perp,p}=\Ker e_{\rm pl}
 \subset C([ -\theta_m,0];\R^{N+1})
\]
in the fold-coordinate history chart, with the maximum history norm, and let
\(\mathscr N_{\perp,p}\) be the corresponding current normal-coordinate
space.  Denote by \(\mathcal A_{\perp,p}(s)\) the derivative of the RFDE in
the normal coordinate furnished by \cref{eq:supp-central-shear}.  Define
\begin{align}
 \mathscr F_{\perp,p}(I)
 &=L^\infty(I;\mathscr N_{\perp,p}),\notag\\
 \mathscr X_{\perp,p}(I)
 &=\left\{\omega\in C(I^-;\mathscr N_{\perp,p}):
       \omega|_I\in W^{1,\infty},\
       \omega_{s_-}\in\mathscr H_{\perp,p}\right\},
 \label{eq:supp-central-normal-spaces}
\end{align}
with graph norm
\[
 \norm\omega_{\mathscr X_{\perp,p}}
 =\sup_{s\in I}\norm{\omega_s}_{\mathscr H_{\perp,p}}
  +\norm{\omega'-\mathcal A_{\perp,p}(\cdot)\omega_\bullet}
          _{\mathscr F_{\perp,p}}.
\]
Then the initial-value operator
\begin{equation}
 \mathbf L_{\perp,p,I}:\mathscr X_{\perp,p}(I)
 \longrightarrow
 \mathscr F_{\perp,p}(I)\times\mathscr H_{\perp,p},
 \qquad
 \omega\longmapsto
 \left(\omega'-\mathcal A_{\perp,p}(\cdot)\omega_\bullet,
       \omega_{s_-}\right)
 \label{eq:supp-central-normal-path-operator}
\end{equation}
is an isomorphism.  There exist \(M_\perp\ge0\) and \(C>0\), independent of
\(N\), such that both \(\mathbf L_{\perp,p,I}^{-1}\) and its first derivative
with respect to \(p\) have norm at most \(C\delta^{-M_\perp}\).
\end{lemma}

\begin{proof}
Since \(e_{\rm pl}\mathcal I_p=\Id\), the inverse of
\cref{eq:supp-central-shear} is
\[
 (u,\omega)\longmapsto\mathcal I_p(u)+\omega.
\]
Subtract the derivative of the graph invariance equation from the full
variational RFDE.  The normal equation is triangular up to a Volterra
term.  Its history-translation component in the fold variables is the
left-translation equation with
an integral over at most one maximal delay, while its transverse present
component has the form
\begin{equation}
 \rho'(s)=\delta^{-1}A_N\rho(s)
          +\mathcal C_p(s)[\kappa_s,\rho_s]+f(s).
 \label{eq:supp-central-normal-equation}
\end{equation}
By the Dobrushin estimate,
\begin{equation}
 \sup_s\norm{\int_0^s
 e^{A_N(s-\zeta)/\delta}f(\zeta)\,\dd\zeta}_N
 \le\frac{\delta}{D\gamma}\norm f_\infty.
 \label{eq:supp-central-poisson}
\end{equation}
Feedback from the normal history into this history-translation component carries
one factor \(\delta\).  On \(I\) its Volterra norm is
\[
 C\delta e^{cS_\delta}\langle S_\delta\rangle^M=o(1).
\]
A Neumann series therefore gives the bounded solution operator, with at most a fixed
power of \(\delta^{-1}\) from the trace norms.  Differentiating the solved
Volterra equation in \(p\) uses the same inverse.  This proves a bound for
solution derivatives; no differentiability assertion for an inverse RFDE
semiflow is involved.
\end{proof}

\subsection{The stable manifold and a truncated boundary-value problem}

Fix \(0<\chi<1\) and \(\sigma_0>\theta_m+1\).  With
\(S_\delta=\sqrt{2p_0\log(1/\delta)}\), set
\begin{equation}
 s_b=\chi S_\delta,\qquad
 s_c=S_\delta-\sigma_0,\qquad
 \mathcal E_b=e^{s_b^2/2},\qquad
 a_{b,\delta}=\frac{\delta}{s_b}\mathcal E_b,
 \label{eq:supp-central-cuts}
\end{equation}
and
\begin{equation}
 b_\delta=S_\delta^{-2}+\sqrt\delta\,S_\delta^{7/2},
 \qquad
 \mathcal A_{bc}=\frac12(s_c^2-s_b^2).
 \label{eq:supp-central-overlap-scale}
\end{equation}
For every prescribed \(q>0\), choose \(p_0\) so that
\begin{equation}
 \delta^{-M}e^{-\mathcal A_{bc}+cS_\delta}
       \langle S_\delta\rangle^M\le C_q\delta^q
 \label{eq:supp-central-action-choice}
\end{equation}
for every finite loss \(M\) occurring below.

\begin{lemma}[Stable-manifold graph on a common transverse section]
\label[lemma]{lem:supp-intermediate-section-graphs}
Assume \(\mathrm{(H_{\rm glob})}\).  Let \(\Sigma_b\) be the transverse section
corresponding to \(s_b\), and let
\(O_{b,p}\) be the intersection of the half-line orbit in
\cref{hyp:supp-outer-manifolds}.  Let \(C_{b,p}^{\rm fin}\) be the
repelling trace used in \cref{prop:finite-comparison}.  There is a
compatible history chart
\[
 \Phi_{b,p}(\psi,u),\qquad \Phi_{b,p}(0,0)=O_{b,p},
\]
whose \(u\)-direction is the oriented unstable tangent vector normal to the
section.  Its domain contains a ball of radius \(Rb_\delta\), with
\(R\) independent of \(N\) and \(\delta\).  With the fixed exponent \(M\)
above,
\begin{equation}
 \norm{\mathbf D_p^1\{\Phi_{b,p}^{-1}(C_{b,p}^{\rm fin})\}}
 \le C\delta^{-M}b_\delta.
 \label{eq:supp-common-section-distance}
\end{equation}

In this chart the stable manifold is
\begin{equation}
 \mathcal S_b^-(p)
 =\{\Phi_{b,p}(\psi,F_{b,p}^-(\psi))\},
 \qquad
 \mathfrak m_{b,p}(\Phi_{b,p}(\psi,u))
 =u-F_{b,p}^-(\psi).
 \label{eq:supp-stable-graph-b}
\end{equation}
The graph is \(C^2\) in the state and \(C^1\) in \(p\), and its
\(C^2\) norm is at most \(C\delta^{-M}\).

Let \(F_{b,p}^{\rm fin}\) be the graph obtained by imposing an affine
boundary functional at the truncated endpoint \(\Sigma_c\).  Then
\begin{equation}
 C_{b,p}^{\rm fin}
 =\Phi_{b,p}(\zeta_{C,p},
                  F_{b,p}^{\rm fin}(\zeta_{C,p})),
 \label{eq:supp-finite-trace-incidence}
\end{equation}
and, for every \(q>0\), after choosing \(p_0\) as in
\cref{eq:supp-central-action-choice},
\begin{equation}
 \norm{\mathbf D_p^1(F_{b,p}^--F_{b,p}^{\rm fin})}_{C^1}
 \le C_q\delta^q.
 \label{eq:supp-stable-finite-comparison}
\end{equation}

If \(k_{b,p}^c\) is the tangent vector in the compatible-history section
obtained by transporting from \(\Sigma_0\) to \(\Sigma_b\) the normalized
tangent vector normal to the section, then
\begin{equation}
 a_{0b,p}:=D\mathfrak m_{b,p}(C_{b,p}^{\rm fin})k_{b,p}^c
 \label{eq:supp-central-terminal-pairing}
\end{equation}
satisfies
\begin{equation}
 c_0a_{b,\delta}\le a_{0b,p}\le C_0a_{b,\delta},
 \qquad
 \abs{D_p\log a_{0b,p}}\le C\delta^{-M}.
 \label{eq:supp-central-terminal-pairing-bound}
\end{equation}
\end{lemma}

\begin{proof}
First continue the orbit from
\cref{hyp:supp-outer-manifolds}(iii) forward until all histories on
\(\Sigma_b\) are formed by the same forward RFDE segment.  The half-line
stable manifold and the implicit-function theorem for transverse
intersections give a local
codimension-one manifold on \(\Sigma_b\).  We now record the mixed
boundary-value formulation because it supplies both the domain of the
graph and the comparison estimate.

At \(r=-\delta s_b/(2\alpha)\), the fold tracking estimate and the slow
expansion used in \cref{eq:supp-negative-residual} give
\[
 w=w_{0,N}(r)+\frac{\delta^2}{2\alpha}
     +O(\delta^3S_\delta^2+\delta^2b_\delta)
\]
for the two reference histories.  Their voltage histories differ by
\(O(b_\delta)\) in \cref{eq:supp-fading-history-norm}.  Applying the
transition map and differentiating its defining equation gives
\cref{eq:supp-common-section-distance}.  Thus both histories lie in the
same compatible chart before either graph equation is solved.

Fix \(T>2\theta_m\) and represent the reference orbit segment from
\(\Sigma_b\) to \(\Sigma_L\) on \([0,T]\).  Choose a strictly increasing
\(C^2\) reference clock
\(\chi_0:[-\theta_m,T]\to\R\) which agrees with fold time on the entry and
terminal history intervals.  Let \(I_{\rm ev}\Subset(0,T)\) be a fixed
interval, and choose \(\rho_{\rm ev}\in C^2([ -\theta_m,T])\) such that
\(\rho_{\rm ev}=0\) on \([ -\theta_m,0]\),
\(\rho_{\rm ev}=1\) on \([T-\theta_m,T]\), and
\(\operatorname{supp}\rho_{\rm ev}'\subset I_{\rm ev}\).  The scalar
\(\beta_{\rm clk}\) varies the
travel time through
\begin{equation}
 \chi_{\beta_{\rm clk}}(s)
 =\chi_0(s)+\beta_{\rm clk}\rho_{\rm ev}(s),
 \qquad \abs{\beta_{\rm clk}}<\beta_0,
 \label{eq:supp-clock-family}
\end{equation}
where \(\beta_0\) is chosen so that every \(\chi_{\beta_{\rm clk}}\) is
increasing.  The subscript distinguishes this clock parameter from the fixed
cubic coefficient \(\beta\) in the RFDE.
For a path \(Z=(Z_v,Z_w)\), define
\begin{align}
 \mathcal H_{\beta_{\rm clk}} Z(s)
 &=\left(\theta\mapsto
 Z_v\!\left[\chi_{\beta_{\rm clk}}^{-1}
       \{\chi_{\beta_{\rm clk}}(s)+\theta\}\right],Z_w(s)\right),\notag\\
 \mathcal E_p(Z,\beta_{\rm clk})
 &=Z'-\chi_{\beta_{\rm clk}}'\mathcal F_p^g
       (\mathcal H_{\beta_{\rm clk}} Z).
 \label{eq:supp-clocked-history-residual}
\end{align}

The path space is
\begin{equation}
 \mathscr D_{bL}^{0;2}
 =\left\{Z\in C([ -\theta_m,T];\R^{N+1}):
 Z|_{[0,T]}\in W^{1,\infty},\quad
 Z|_{[-\theta_m,0]\cup I_{\rm ev}\cup[T-\theta_m,T]}
       \in W^{2,\infty}\right\}.
 \label{eq:supp-mixed-path-space}
\end{equation}
Let \(\mathscr Y_b=C([ -\theta_m,0];\R^N)\times\R\) be the entry-history
space in the scaled fold chart.  To specify the source norm, put
\[
 \widehat r_p(s)=r_p(\chi_0(s)),\qquad
 \widehat a_p(s)=a_p(\chi_0(s)),\qquad
 q_p(s)=\{\delta\chi_0'(s)\}^{-1}\widehat r_p'(s)
\]
along the negative reference branch, expressed in the computational clock on
\([0,T]\).  Let
\(\mathscr F_{bL,p}=L^\infty([0,T];\R^{N+1})\), and for
\(F=(F_v,F_w)\in\mathscr F_{bL,p}\) set
\begin{align}
 (\mathcal V_pF_w)(s)
 &=q_p(s)\int_0^s\frac{F_w(\zeta)}{q_p(\zeta)}\,\dd\zeta,\notag\\
 \norm F_{\mathscr F_{bL,p}}
 &=\max\left\{
   \operatorname*{ess\,sup}_{0\le s\le T}
        \frac{\delta}{\widehat a_p(s)}\norm{F_v(s)}_N,\
   \norm{\mathcal V_pF_w}_\infty
   +\operatorname*{ess\,sup}_{0\le s\le T}
        \frac{\delta}{\widehat a_p(s)}\abs{F_w(s)}\right\}.
 \label{eq:supp-mixed-source-norm}
\end{align}
Here the second component is the variation-of-constants norm for the scalar
recovery equation after removal of the time-tangent direction.  Equip
\(\mathscr D_{bL}^{0;2}\) with the graph norm formed from the supremum of
\(\norm{\mathcal H_0Z(s)}_{p,\chi_0(s),\sharp}\),
\(\norm{Z'}_{\mathscr F_{bL,p}}\), and the displayed
\(W^{2,\infty}\) window norms.

Extend \(\mathfrak m_{L,p}\) from \(\mathscr H_{\Sigma_L}^2\) to a fixed
tubular neighborhood by making it constant in the flow-normal coordinate.
For fixed \(\psi\), define
\begin{equation}
 \mathcal R_{bL,p}(Z,\beta_{\rm clk},u;\psi)
 =\left(
 \mathcal E_p(Z,\beta_{\rm clk}),\
 \mathcal H_{\beta_{\rm clk}} Z(0)-\Phi_{b,p}(\psi,u),\
 r(\mathcal H_{\beta_{\rm clk}} Z(T))+r_L,\
 \mathfrak m_{L,p}(\mathcal H_{\beta_{\rm clk}} Z(T))
 \right).
 \label{eq:supp-mixed-residual}
\end{equation}
This defines a map
\begin{equation}
 \mathcal R_{bL,p}:
 \underbrace{\mathscr D_{bL}^{0;2}\times\R_{\beta_{\rm clk}}\times\R_u}
             _{\mathbb X_{bL,p}}
 \longrightarrow
 \underbrace{\mathscr F_{bL,p}\times\mathscr Y_b
             \times\R_{\rm ev}\times\R_{\rm st}}
             _{\mathbb Y_{bL,p}}.
 \label{eq:supp-mixed-residual-spaces}
\end{equation}
The entry row is equality in the compatible history space.  The
terminal-section row
places the terminal history on \(\Sigma_L\), and the last row imposes the
stable-manifold condition there.  If
\((\bar Z_p,0,\bar u_p;\bar\psi_p)\) is the reference zero corresponding to
the stable-manifold orbit, set
\begin{equation}
 \mathbb B_p
 =D_{(Z,\beta_{\rm clk},u)}\mathcal R_{bL,p}
       (\bar Z_p,0,\bar u_p;\bar\psi_p)
 :\mathbb X_{bL,p}\longrightarrow\mathbb Y_{bL,p}.
 \label{eq:supp-mixed-linearization}
\end{equation}

The quotient variational equation has an exponential dichotomy with a
codimension-one stable bundle and a one-dimensional unstable bundle.
Stable data are propagated forward.  Only the unstable scalar is
determined from the terminal boundary functional.  To see the latter
explicitly, let
\(\widehat y_0\) be the terminal history of a particular solution and
\(\widehat h^u\) the terminal history of the homogeneous unstable solution.
For prescribed terminal value \(b\), the replacement
\begin{equation}
 \widehat y=\widehat y_0+\widehat h^u
 \frac{b-D\mathfrak m_{L,p}(\bar Z_{-,L})\widehat y_0}
      {D\mathfrak m_{L,p}(\bar Z_{-,L})\widehat h^u}
 \label{eq:supp-terminal-schur}
\end{equation}
solves the final row.  All three objects in
\cref{eq:supp-terminal-schur} lie in the same compatible terminal-history
tangent--dual pairing.  By \cref{eq:supp-terminal-normalization}, and by
continuity on the stated neighborhood, its denominator is at least
\(3/4\) for small \(\delta\).

The exponential-dichotomy Green operator, the scalar recovery
variation-of-constants formula, and the terminal-section row therefore make the
linearization of \cref{eq:supp-mixed-residual} an isomorphism.  In the
norms weighted by the integrated rate \(\mathcal A_u\), its inverse is bounded by
\(C\log(1/\delta)\).  On the radius-\(Rb_\delta\) ball, Taylor's formula
gives
\[
 \sup_X\norm{\mathbb B_p^{-1}
       \{D\mathcal R_{bL,p}(X)-\mathbb B_p\}}
 \le CR\log(1/\delta)b_\delta<\frac13
\]
after \(R\) is fixed and \(\delta_0\) is decreased.  The
parameter-dependent contraction theorem proves
\cref{eq:supp-stable-graph-b}, its stated domain, and its derivatives.

For the truncated boundary-value problem, replace the last component in
\cref{eq:supp-mixed-residual} by the affine row at \(\Sigma_c\) that
vanishes on \(C_{b,p}^{\rm fin}\).  The two problems have the same
equation, entry trace, phase row, and stable coordinate between
\(\Sigma_b\) and \(\Sigma_c\).  They differ only in the terminal unstable
scalar.  Propagating that scalar backward on the one-dimensional unstable
bundle yields
\[
 \norm{\mathbf D_p^1(F_{b,p}^--F_{b,p}^{\rm fin})}_{C^1}
 \le C\delta^{-M}e^{-\mathcal A_{bc}+cS_\delta}
       \langle S_\delta\rangle^M.
\]
\Cref{eq:supp-central-action-choice,eq:supp-finite-trace-incidence}
give \cref{eq:supp-stable-finite-comparison}.

It remains to evaluate the section-normal tangent vector at \(\Sigma_0\).  For
the singular fold equation
the independent solution normalized by the first integral is
\[
 n(s)=\binom{I(s)}{\mathcal E(s)-sI(s)},\qquad
 \mathcal E(s)=e^{s^2/2},\qquad
 I(s)=\int_0^s\mathcal E(\zeta)\,\dd\zeta.
\]
Intersecting with a fixed-\(r\) transverse section removes the time component
\(I(s)\) and leaves
\(\mathcal E(s)\) in the section-normal direction.  In the variables of
\cref{eq:rw-fold-coordinates}, the recovery component is
\(-\delta^2\mathcal E(s)\{1+o(1)\}\).  Division by
\[
 w_{0,N}'(-\delta s_b/(2\alpha))
 =\delta s_b\{1+o(1)\}
\]
after normalization by the recovery-coordinate derivative gives
\(a_{0b,p}\asymp\delta\mathcal E_b/s_b\).  The fold graph bounds, the
boundary-value inverse in \cref{lem:supp-central-boundary-inverse}, and the
\(C^1\) transition map give the logarithmic derivative estimate in
\cref{eq:supp-central-terminal-pairing-bound}.
\end{proof}

\subsection{Intersection on the fold section}

\begin{lemma}[Unique local intersection with the stable manifold]
\label[lemma]{lem:supp-stable-manifold-intersection}
Assume \(\mathrm{(H_{\rm glob})}\).  Let \(\mathfrak C_{0,p}\) be the chart in
\cref{eq:connection-central-chart}.  Define the stable-manifold membership
function pulled back from \(\Sigma_b\) by
\begin{equation}
 \mathcal M_{0,p}(\psi,u)
 =\mathfrak m_{b,p}\!\left(
   \mathcal P_{\Sigma_0\Sigma_b,p}
      \mathfrak C_{0,p}(\psi,u)\right).
 \label{eq:supp-pulled-membership}
\end{equation}
Write the finite repelling trace as
\(\mathfrak C_{0,p}(\bar\psi_p,\bar u_p)\), and introduce
\begin{equation}
 \xi=a_{0b,p}(u-\bar u_p).
 \label{eq:supp-scaled-unstable-coordinate}
\end{equation}
For every prescribed \(Q>12\), the exponent \(p_0\) in
\cref{eq:supp-central-cuts} can be fixed so that, on
\begin{equation}
 \norm{\psi-\bar\psi_p}
 \le c_r\delta^{Q+M+2},
 \qquad
 \abs{\xi}\le C_r\delta^Q,
 \label{eq:supp-intersection-rectangle}
\end{equation}
one has
\begin{equation}
 \partial_\xi
 \mathcal M_{0,p}\!\left(
 \psi,\bar u_p+a_{0b,p}^{-1}\xi\right)\ge\frac12.
 \label{eq:supp-membership-slope}
\end{equation}
The two \(\xi\)-faces have opposite signs.  Hence the zero set is a unique
\(C^1\) graph
\[
 u=F_{N,\delta,p}^-(\psi)
\]
on the indicated \(\psi\)-ball.  Every history in this graph, together with
every translated history segment required by the delay terms, remains in the
stated fold neighborhood.
\end{lemma}

\begin{proof}
By \cref{eq:supp-finite-trace-incidence,%
eq:supp-stable-finite-comparison}, applied with exponent \(Q+M+4\),
\begin{equation}
 \norm{\mathbf D_p^1
   \mathcal M_{0,p}(\bar\psi_p,\bar u_p)}
 \le C\delta^{Q+M+4}.
 \label{eq:supp-reference-membership}
\end{equation}
Put
\[
 H_p(\psi,\xi)=
 \mathcal M_{0,p}
   (\psi,\bar u_p+a_{0b,p}^{-1}\xi).
\]
The fold graph estimates, the boundary-value inverse in
\cref{lem:supp-central-boundary-inverse}, and
\cref{eq:supp-central-terminal-pairing-bound} imply
\begin{equation}
 \norm{D_\psi H_p}\le C\delta^{-M},
 \qquad
 \abs{\partial_\xi H_p-1}\le\frac12,
 \qquad
 \norm{D_{(\psi,\xi)}\partial_\xi H_p}
 \le C\delta^{-M}
 \label{eq:supp-rectangle-derivatives}
\end{equation}
on \cref{eq:supp-intersection-rectangle}.

The state displacement produced by \(\xi\) is bounded by
\(C(s_b/\delta)\abs{\xi}\), and the displacement produced by
\(\psi-\bar\psi_p\) is bounded by
\(C\delta^{-M}\norm{\psi-\bar\psi_p}\).  Since \(Q>12\), both are
smaller than the chart radius
\(Rb_\delta\).  The same estimates applied to all translated history
segments required by the delay terms keep the orbit inside the fold
neighborhood and preserve the transversality margin.

Choose \(c_r\) small and then \(C_r\) larger than the constant in
\cref{eq:supp-reference-membership}.  The \(\psi\)-variation is less than
one quarter of the face height.  Integrating
\cref{eq:supp-rectangle-derivatives} in \(\xi\) gives opposite signs on
the two faces, while \cref{eq:supp-membership-slope} gives uniqueness.
The parameter-dependent implicit-function theorem now yields the asserted
graph.
\end{proof}

\subsection{Comparison with the finite-interval function}

\begin{lemma}[Comparison of the two defining functions]
\label[lemma]{lem:supp-defining-function-comparison}
Assume \(\mathrm{(H_{\rm glob})}\).  Let \(G_{N,\delta}^{g}\) be the
stable-manifold defining function in
\cref{eq:connection-defining-function}, and let \(D_N^{\rm fin}\) be the
finite-interval function in \cref{prop:finite-comparison}.  For every
prescribed \(Q>12\), after choosing \(p_0\) in
\cref{eq:supp-central-action-choice} with exponent larger than \(Q+M+4\),
\begin{equation}
 \norm{\mathbf D_p^1
       \{G_{N,\delta}^{g}-D_N^{\rm fin}\}}
 \le C\delta^Q.
 \label{eq:supp-defining-function-comparison}
\end{equation}
This estimate compares a defining function for the stable-manifold
connection of one fixed RFDE with a finite-interval calculation.  It does
not identify the zero set of the latter with a heteroclinic connection.
\end{lemma}

\begin{proof}
Use the chart normalization in
\cref{eq:connection-chart-normalization}.  The attracting and repelling
finite traces have coordinates
\[
 A_p^{\rm fin}=\mathfrak C_{0,p}(\bar\psi_p,u_p^a),
 \qquad
 C_p^{\rm fin}=\mathfrak C_{0,p}(\bar\psi_p,u_p^r),
\]
and the finite-interval construction gives
\begin{equation}
 D_N^{\rm fin}=u_p^a-u_p^r.
 \label{eq:supp-finite-coordinate-difference}
\end{equation}
Write the incoming intersection point as
\[
 A_p^0=\mathfrak C_{0,p}(\psi_A,u_A).
\]
By \cref{eq:supp-positive-history-comparison}, applied with exponent
\(Q+M+4\) to absorb the polynomial chart bounds,
\begin{equation}
 \norm{\mathbf D_p^1
   \{(\psi_A,u_A)-(\bar\psi_p,u_p^a)\}}
 \le C\delta^{Q+M+4}.
 \label{eq:supp-positive-coordinate-comparison}
\end{equation}

On the negative side,
\cref{lem:supp-intermediate-section-graphs,%
lem:supp-stable-manifold-intersection} compare the stable manifold supplied by
\(\mathrm{(H_{\rm glob})}\) and the graph defined by the truncated boundary
condition on a common forward orbit segment.  Pulling both
membership functions back to \(\Sigma_0\), without inverting a history-space
transition map, gives
\begin{equation}
 \norm{\mathbf D_p^1
   \{F_{N,\delta,p}^-(\bar\psi_p)-u_p^r\}}
 \le C\delta^{Q+4}.
 \label{eq:supp-negative-coordinate-comparison}
\end{equation}
The state and parameter derivatives of \(F_{N,\delta,p}^-\) may have a
fixed polynomial loss.  The extra margin in
\cref{eq:supp-positive-coordinate-comparison} absorbs that loss when the
mean-value theorem is applied between \(\psi_A\) and \(\bar\psi_p\).

Finally,
\[
 G_{N,\delta}^{g}-D_N^{\rm fin}
 =(u_A-u_p^a)
  -\{F_{N,\delta,p}^-(\psi_A)-u_p^r\}
\]
by \cref{eq:connection-defining-function,%
eq:supp-finite-coordinate-difference}.  Insert
\cref{eq:supp-positive-coordinate-comparison,%
eq:supp-negative-coordinate-comparison} and differentiate the already
constructed \(C^1\) families.  This proves
\cref{eq:supp-defining-function-comparison}.
\end{proof}
